\labsection{9}{Validate an imbalanced classifier without leakage}{sec:lab09}

\begin{labproject}{Validation and preprocessing: evaluate the process, not one lucky split}
\textbf{Coverage.} Chapter~\ref{ch:validation}.\\
\textbf{Goal.} You should be able to construct stratified folds, prevent preprocessing leakage, compare thresholds across folds, quantify metric variability, and distinguish discrimination from calibration.\\
\textbf{Datasets.} \texttt{credit\_card\_fraud\_dataset.csv}; Heart Failure is used for a calibration extension.\\
\textbf{Scenario.} Fraud is rare, so a superficially high accuracy can be operationally useless. You must evaluate a complete modelling process with stratified cross-validation, fold-local preprocessing, and threshold-sensitive metrics.

\courselabflow{Create folds}{Fit inside fold}{Compare costs}{Check calibration}


\textbf{Run the complete script.} Start a clean R session in the course folder. Each step continues the same file; comment headings identify the operation performed by the next block.
\begin{lstlisting}[style=utpterminal]
Rscript all_experiments/lab09_validation.R
\end{lstlisting}

\noindent\textbf{Step 1.} Load, validate and prepare the data
\lstinputlisting[style=utpr,firstline=1,lastline=11]{all_experiments/lab09_validation.R}

\noindent\textbf{Step 2.} Create stratified fold labels.
\courselabtrace{Repeat preparation inside each fold}{Fit}{Use the fold-training subset to learn parameters}{Transform}{Apply those parameters to the validation subset}{Aggregate}{Combine validation metrics; never train on held-out rows}

\lstinputlisting[style=utpr,firstline=12,lastline=43]{all_experiments/lab09_validation.R}


\textbf{Use the chapter while coding.} Revisit \hyperref[ch:validation]{the earlier chapter explanation}; match its operation to the variables and state changes below.

\textbf{Follow the supplied implementation.} Inspect the scaling assignments inside the loop, not just the final scores. The cost\_5\_to\_1 summary uses an explicitly assumed cost ratio. Compare both mean and variability before selecting a decision rule.
\courselabtrace{Trace preprocessing inside a validation fold}{1. Input / state}{K = 5. fold selects the current validation subset and its complementary training rows.}{2. Operation}{Estimate the scaling statistics on fold-training rows, then reuse them for validation.}{3. Check}{Collect results at thresholds .10, .25 and .50; aggregate across the same folds.}


\lstinputlisting[style=utpr,firstline=44,lastline=75]{all_experiments/lab09_validation.R}
\lstinputlisting[style=utpr,firstline=76,lastline=81]{all_experiments/lab09_validation.R}

\noindent\textbf{Step 3.} Illustrative operational choice: missing a fraud costs 5 times a false alarm.
\lstinputlisting[style=utpr,firstline=82,lastline=88]{all_experiments/lab09_validation.R}

\noindent\textbf{Step 4.} Calibration extension with Heart Failure probabilities on a fresh holdout split.
\lstinputlisting[style=utpr,firstline=89,lastline=113]{all_experiments/lab09_validation.R}
\textbf{Chapter connection.} Revisit \hyperref[ch:validation]{the chapter explanation} and map its states to the steps below.
% Reused from chapters/09-validation-preprocessing.tex
\begin{labvisual}{Recall the chapter model}
\begin{labgraphic}
\begin{tikzpicture}[x=\linewidth,y=1cm,
  foldcell/.style={draw=UTPBlue!70,fill=UTPPaleBlue,rounded corners=1mm,
    line width=0.6pt,text width=0.13\linewidth,minimum height=0.62cm,
    inner xsep=1mm,inner ysep=1mm,align=center,
    font=\scriptsize\color{UTPNavy}},
  heldout/.style={foldcell,draw=UTPTeal!80,fill=UTPPaleTeal,
    font=\scriptsize\bfseries\color{UTPNavy}},
  heading/.style={font=\scriptsize\bfseries\color{UTPNavy},inner sep=0pt}]
\node[heading,anchor=east] at (0.095,0.65) {Run};
\foreach \c in {0,...,4}{
  \pgfmathtruncatemacro{\foldnumber}{\c+1}
  \node[heading] at (0.20+\c*0.16,0.65) {Fold \foldnumber};
}
\foreach \r in {0,...,4}{
  \pgfmathtruncatemacro{\runnumber}{\r+1}
  \node[heading,anchor=east] at (0.095,-\r*0.78) {\runnumber};
  \foreach \c in {0,...,4}{
    \ifnum\r=\c
      \node[heldout] at (0.20+\c*0.16,-\r*0.78) {Validate};
    \else
      \node[foldcell] at (0.20+\c*0.16,-\r*0.78) {Train};
    \fi
  }
}
\node[font=\scriptsize\color{UTPMuted},align=center,
  text width=0.86\linewidth] at (0.50,-3.85)
  {Each run uses four training folds and one held-out validation fold.};
\end{tikzpicture}
\end{labgraphic}\end{labvisual}


\begin{verification}
\textbf{Explain and extend your result.} Compare two thresholds at the stated cost ratio, then repeat with a second ratio. Explain whether the best decision threshold changes even when predicted probabilities stay fixed.


\textbf{Required evidence.} Submit the majority-class baseline, five-fold metric summary, one variability result, and a threshold recommendation tied to an explicit error-cost assumption. Include one sentence identifying where leakage would occur if global scaling were applied before cross-validation.

\textbf{Professional standard.} Validation must reproduce the full future prediction process. Any transformation learned from all records before validation contaminates the estimate of generalisation performance.
\end{verification}
\end{labproject}
